% =====================================================================
% Paper E1 v69 (research-article rewrite of paperE1_cod_forecast_ladder_v59.tex;
% repo v60 on branches e2-v3-source-year / lean-audit-v4 = v59 + Code-availability subsection + extended AI declaration, both carried here).
% Numbered v69 because branch edwards-framework-e1-audit-implementation holds a parallel line e1/paperE1_cod_forecast_ladder_v49..v68 (not merged here);
% version numbers are kept unique across all branches.
% Does a surplus-production ladder improve forecasts of Northern cod?
% A scored retention test against persistence on NAFO 2J3KL
% Author: Amin Abaee (Independent Researcher, Tehran, Iran)
% Target: Fisheries Research (elsarticle, author-year)
%
% v69 (2026-10-15): restructured as a research article. Main text carries the
% design, the pooled verdict, a new period decomposition (Table 6, Fig. 4),
% the identification argument (Proposition 1, Fig. 5), the operating
% characteristics of the rule (Table 8, Fig. 6) and the 2016+ verdict layer.
% Variant passes, the full Diebold-Mariano / bootstrap grids, fitted
% parameters, the floor-forecast census, secondary scores, the simulation
% design, the suspension chronology, the verification record and the Edwards
% Aquifer restatement are in the Supplementary Information. No archived
% number was changed; all four runners were re-executed in a third software
% environment before this revision (SI S13). Compiles with pdfLaTeX /
% XeLaTeX / LuaLaTeX / tectonic.
% =====================================================================
\documentclass[preprint,11pt,authoryear]{elsarticle}

\usepackage{amsmath,amssymb,amsthm}
\usepackage{booktabs}
\usepackage{array}
\usepackage{tabularx}
\usepackage{multirow}
\usepackage{graphicx}
\graphicspath{{./}{figs/}{paperE1_cod_forecast_ladder_v69_src/figs/}}
\usepackage{enumitem}
\usepackage{placeins}
\renewcommand{\topfraction}{0.9}
\renewcommand{\bottomfraction}{0.7}
\renewcommand{\textfraction}{0.07}
\renewcommand{\floatpagefraction}{0.75}
\usepackage[colorlinks=true,urlcolor=blue,linkcolor=blue,citecolor=blue]{hyperref}
\emergencystretch=3em
\setlength{\tolerance}{2000}

\newtheorem{proposition}{Proposition}
\theoremstyle{definition}
\newtheorem{definition}{Definition}
\theoremstyle{remark}
\newtheorem{remark}{Remark}

\newcommand{\kt}{\,\mathrm{kt}}
\newcommand{\Blim}{B_{\mathrm{lim}}}
\newcommand{\Spos}[1]{\left[#1\right]_{+}}
\newcommand{\RMSE}{\mathrm{RMSE}}
\newcommand{\MAE}{\mathrm{MAE}}
\newcommand{\pers}{\mathrm{persist}}
\newcommand{\sA}{\mathfrak{s}}
\newcolumntype{Y}{>{\centering\arraybackslash}X}
\newcolumntype{L}{>{\raggedright\arraybackslash}X}

\journal{Fisheries Research}

\begin{document}

\begin{frontmatter}

\title{Does a surplus-production ladder improve forecasts of Northern cod? A scored retention test against persistence on NAFO 2J3KL}

\author[1]{Amin Abaee\corref{cor1}}
\ead{amin_abaee@ut.ac.ir}
\cortext[cor1]{Corresponding author. ORCID: 0000-0002-0019-1842.}
\affiliation[1]{organization={Independent Researcher}, city={Tehran}, country={Iran}}

\begin{abstract}
Prediction-skill diagnostics in stock assessment are usually computed on the fitted index, whereas advice concerns the biomass state. We score a five-module surplus-production ladder (autonomous, depensatory, catch-driven stock-flow, autoregressive-residual and lagged-start modules) against last-value persistence on the spawning-stock biomass of Northern cod (\emph{Gadus morhua}, NAFO 2J3KL) under a retention rule whose scoring core was coded before the first scoring pass: a module is kept only if it lowers rolling-origin RMSE against both persistence and the next-simpler module, at one- and five-year horizons, by more than a 5\% band. On the 1983--2015 NCAM series and, separately, the 1954--2024 xteNCAM series, no module is retained: one-year RMSE is 98\,kt for persistence against 115--196\,kt for the modules, and 84 against 119--206\,kt; paired moving-block bootstrap intervals contain zero on the first series and exclude it on the second. A period decomposition, added here, shows that the pooled deficit is carried by the 1991--1995 collapse, that the autonomous modules forecast the numerical floor at every post-collapse origin because training-mean catch exceeds production, and that catch-supplied modules have one-step skill in the low-biomass regime on one specification only. Under known truth the rule retains a true autonomous module in 97\% of replicates but a true stock-flow or depensation module in under 15\%, and over-retains when the truth lies outside the class; non-retention is therefore strong evidence against the autonomous module and weak evidence against the others. The collapse is missed because the lower equilibrium of a catch-driven Schaefer map is generated by removals, not by depensation. After 2015 the reference point and assessment vintage moved more than forecast skill did.
\end{abstract}

\begin{keyword}
Northern cod \sep biomass forecasting \sep surplus production \sep persistence benchmark \sep rolling-origin evaluation \sep Diebold--Mariano test \sep model retention \sep limit reference point
\end{keyword}

\end{frontmatter}
